% Full extended-mean, renewal, empirical-CDF and finite-alphabet arguments.
\section{Applications of strong laws: proofs}
\label{sec:annex-p15-strong-law-applications}

The statements are in Chapter~\ref{sec:p15-strong-law-applications-chapter}.
The source arguments are Durrett's Theorem 2.4.5, the full renewal
application, the random-index observation, the discontinuous-CDF proof
of Theorem 2.4.9 and the information-sum proof in Example 2.4.10
\cite[pp.~78--80]{DurrettPTEAuthor2019}.

\begin{annexitem}{annex:p15.extended-mean-strong-law}{thm:p15-extended-mean-strong-law}{Clipping from above, both signs and the common event}
\begin{proof}
If both part expectations are finite, then \(\mathbb E|X_1|<\infty\),
and Theorem~\ref{thm:p15-etemadi-strong-law} gives the finite
limit. Suppose instead that \(\mathbb EX_1^+=\infty\) and
\(\mathbb EX_1^-<\infty\). For each positive integer \(M\), define
\[
 X_i^{[M]}=X_i\wedge M=\min(X_i,M),\qquad
 S_n^{[M]}=\sum_{i=1}^nX_i^{[M]}.
\]
These variables are i.i.d. measurable transforms of the originals.
Their positive parts are \(X_i^+\wedge M\) and their negative
parts are exactly \(X_i^-\). Hence
\(\mathbb E|X_i^{[M]}|\leq M+\mathbb EX_1^-<\infty\).
The same integrable strong law yields
\(S_n^{[M]}/n\to\mathbb EX_1^{[M]}\) almost surely for each \(M\).
Intersect these events over the positive integers. On the resulting
probability-one event, \(S_n\geq S_n^{[M]}\) gives
\[
 \liminf_{n\to\infty}S_n/n\geq\mathbb EX_1^{[M]}
             =\mathbb E(X_1^+\wedge M)-\mathbb EX_1^-
                         \quad\text{for every }M\geq1.
\]
Monotone convergence makes the first expectation increase to infinity;
the subtracted expectation is finite and fixed. The lower bound
therefore tends to infinity, proving \(S_n/n\to+\infty\).
Only a countable family of clipping levels is needed to obtain one
common event.

If \(\mathbb EX_1^-=\infty\) and \(\mathbb EX_1^+<\infty\),
apply the preceding argument to \(-X_i\), whose positive and negative
parts are interchanged. It gives \(-S_n/n\to+\infty\), which is
the asserted negative limit. These cases exhaust the domain on which
the extended expectation is defined. No subtraction of two infinite
expectations has been made.
\end{proof}
\end{annexitem}

\begin{annexitem}{annex:p15.random-index-limits}{thm:p15-random-index-limits}{Pointwise subsequences and the complete counterexample family}
\begin{proof}
For any Borel set \(B\) in the extended real line,
\[
 \{Z_{J_m}\in B\}
   =\bigcup_{j\geq1}\bigl(\{J_m=j\}\cap\{Z_j\in B\}\bigr),
\]
so the sampled variable is measurable. Intersect the two
probability-one convergence events. At an outcome in their intersection,
the deterministic sequence \(Z_j(\omega)\) tends to \(Z(\omega)\)
and the integer sequence \(J_m(\omega)\) tends to infinity.
Every sufficiently late sampled index therefore lies beyond the index
required by the definition of convergence of \(Z_j(\omega)\).
This proves the conclusion, for finite limits and for either infinite
limit, without an independence premise.

For the counterexample, the Bernoulli laws give
\(\mathbb P(|B_n|>\varepsilon)=a_n\to0\) for \(0<\varepsilon<1\);
for \(\varepsilon\geq1\) this probability is zero. Thus \(B_n\to0\)
in probability. The independent events \(\{B_n=1\}\) have a divergent
probability sum. The second Borel--Cantelli lemma gives an event
\(H\) of probability one on which infinitely many ones occur.
For a fixed \(m\), the extended first-hit index
\(K_m=\inf\{j\geq m:B_j=1\}\), with \(\inf\varnothing=\infty\),
is measurable: for finite \(j\geq m\), its event \(\{K_m=j\}\)
requires \(B_j=1\) and all intervening \(B_i=0\); the event
\(\{K_m=\infty\}\) is the complement of
\(\bigcup_{j\geq m}\{K_m=j\}\) and is measurable as well.
Replacing an infinite \(K_m\) by \(m\) therefore defines the stated
finite measurable \(J_m\), with \(J_m\geq m\) at every outcome.
On \(H\), every \(K_m\) is finite and \(B_{J_m}=1\) for every
\(m\). The proposed analogue for convergence in probability consequently
fails. The example \(a_n=1/(n+1)\) satisfies the conditions because
its terms tend to zero and its harmonic sum diverges.
\end{proof}
\end{annexitem}

\begin{annexitem}{annex:p15.renewal-inverse-law}{thm:p15-renewal-inverse-law}{The last tied epoch and the finite or infinite inverse limit}
\begin{proof}
The events \(\{X_1\geq1/j\}\), \(j\geq1\), increase in union
to \(\{X_1>0\}\). If the latter has positive probability, at least
one of the former does. For that \(j\),
\(\mathbb EX_1\geq j^{-1}\mathbb P(X_1\geq1/j)>0\).
Thus the extended mean \(\mu\) lies in \((0,\infty]\).
Nonnegativity makes its negative-part expectation zero, so the
extended-mean strong law applies and gives \(T_n/n\to\mu\)
on a measurable event \(\Omega_0\) of probability one.
At any outcome of this event, \(T_n\to\infty\): for a finite
positive mean the ratios are eventually at least \(\mu/2\),
and for an infinite mean they are eventually at least one.

Fix such an outcome. The sequence \(T_n\) is nondecreasing, every
term is finite, and it tends to infinity. For each finite \(t\geq0\),
the set of indices with \(T_n\leq t\) is nonempty, since it includes
zero, and is finite, since sufficiently late \(T_n\) exceed \(t\).
It therefore has a largest element \(N(t)\), and
\[
 T_{N(t)}\leq t<T_{N(t)+1}.
\]
The second inequality is strict by maximality, even when some earlier
epochs coincide. All zeros in any finite tied group have been included
in the count. For every integer \(q\), \(t\geq T_q\) implies
\(N(t)\geq q\). Thus \(N(t)\to\infty\) as real \(t\to\infty\).
In particular it is positive for all sufficiently large \(t\).

For those large times, divide the epoch inequalities by \(N(t)\):
\[
 \frac{T_{N(t)}}{N(t)}
   \leq\frac{t}{N(t)}
   <\frac{T_{N(t)+1}}{N(t)+1}\,
                       \frac{N(t)+1}{N(t)}.
\]
At this fixed outcome, the deterministic limit \(T_n/n\to\mu\)
holds along every index tending to infinity. It therefore holds along
both \(N(t)\) and \(N(t)+1\), simultaneously for unrestricted
real \(t\to\infty\); also \((N(t)+1)/N(t)\to1\).
If \(\mu<\infty\), the two bounds tend to \(\mu>0\),
so \(t/N(t)\to\mu\) and \(N(t)/t\to1/\mu\).
If \(\mu=\infty\), the lower bound alone implies
\(t/N(t)\to\infty\), hence \(N(t)/t\to0\).
Every step holds on the same measurable \(\Omega_0\), rather than
an intersection over uncountably many times.

For fixed \(t\geq0\) and integer \(q\geq1\), monotonicity of the
epochs gives
\(\{N(t)\geq q\}=\{T_q\leq t\}\), including outcomes with
\(N(t)=\infty\). These measurable threshold events establish
measurability of the extended integer count.
If \(\mathbb P(X_1>0)=0\), each nonnegative \(X_i\) equals zero
almost surely. Intersect these countably many events. On their
intersection every \(T_n=0\), so every nonnegative \(t\) admits
all indices and \(N(t)=\infty\). This proves the separate degenerate
case as well.
\end{proof}
\end{annexitem}

\begin{annexitem}{annex:p15.empirical-cdf-structure}{thm:p15-empirical-cdf-structure}{Right limits, countable suprema and grids through jumps}
\begin{proof}
At a fixed outcome, \(\nu_n\) is the finite average of the probability
measures concentrated at \(X_1,\ldots,X_n\). It is nonnegative,
countably additive and has total mass one; countable additivity follows
by interchanging a finite sum and a nonnegative countable sum.
Its CDF and left limits are therefore those stated, by
Theorem~\ref{thm:p15-cdf-laws}. For fixed real \(x\), both \(F_n(x)\)
and \(F_n(x-)\) are measurable random variables, since they are
finite averages of event indicators.

For every real \(x\), choose rational numbers decreasing to \(x\)
from above. Right continuity of both CDFs gives convergence of their
difference along these rationals to \(F_n(x)-F(x)\), at each
outcome. Thus \(|F_n(x)-F(x)|\) is at most the supremum of the
absolute differences at rational arguments. The reverse bound follows
by inclusion of the rational arguments among the real ones. This proves
the countable-supremum identity for \(D_n\); consequently \(D_n\)
is measurable and lies in \([0,1]\).

Fix a level \(u=j/k\) strictly between zero and one. The CDF tail
limits ensure that \(\{y:F(y)\geq u\}\) is nonempty and bounded
below, so its infimum \(x_{j,k}\) is finite. If \(z<x_{j,k}\),
then \(F(z)<u\); taking \(z\uparrow x_{j,k}\) yields
\(F(x_{j,k}-)\leq u\). For every \(h>0\), there is a member
of the defining set below \(x_{j,k}+h\), so monotonicity gives
\(F(x_{j,k}+h)\geq u\). Right continuity as \(h\downarrow0\)
yields \(F(x_{j,k})\geq u\). Increasing the level shrinks its
upper-level set and cannot decrease its infimum, proving the grid's
monotonicity.

For an interior adjacent pair,
\[
 F(x_{j,k}-)-F(x_{j-1,k})
           \leq j/k-(j-1)/k=1/k.
\]
At \(j=1\) the lower function value is zero, and at \(j=k\)
the upper left-limit value is one, so the same estimate follows
from the corresponding one interior inequality. If two adjacent
points coincide, the left side is nonpositive; the estimate still
holds and their intervening open interval is empty.
\end{proof}
\end{annexitem}

\begin{annexitem}{annex:p15.continuous-cdf-uniform}{thm:p15-continuous-cdf-uniform}{A finite value grid in the continuous case}
\begin{proof}
Fix an integer \(k\geq2\) and use the quantile grid of \(F\).
Continuity makes \(F(x_{j,k}-)=F(x_{j,k})\), so the two
quantile inequalities force \(F(x_{j,k})=j/k\).
Distinct interior levels therefore give distinct grid points.
Pointwise convergence at the finitely many interior grid points
gives an \(N_k\) after which
\(|G_n(x_{j,k})-F(x_{j,k})|\leq1/k\) for every such \(j\).
Give both CDFs their values zero and one at the formal infinite
endpoints, whose errors are identically zero.

For a real \(x\) between two adjacent grid points, monotonicity
bounds \(G_n(x)\) between its two endpoint values, and bounds
\(F(x)\) between the corresponding endpoint values of \(F\).
The latter values differ by \(1/k\). Hence
\(|G_n(x)-F(x)|\leq2/k\) for \(n\geq N_k\).
At an interior grid point the error is already at most \(1/k\).
The first and last intervals are also covered, because CDFs take
values between zero and one and have the assigned endpoint values.
Thus the supremum is at most \(2/k\) after \(N_k\).
Letting \(k\to\infty\) proves uniform convergence.
This uses the tail normalization of \(F\) to make the finite-level
grid reach the full range from zero to one.
\end{proof}
\end{annexitem}

\begin{annexitem}{annex:p15.glivenko-cantelli}{thm:p15-glivenko-cantelli}{The full grid proof for discontinuous distributions}
\begin{proof}
For a fixed real \(x\), the i.i.d. bounded variables
\(\mathbf1_{\{X_i\leq x\}}\) have mean \(F(x)\), so the
integrable strong law, Theorem~\ref{thm:p15-etemadi-strong-law},
gives \(F_n(x)\to F(x)\) almost surely.
Apply it separately to \(\mathbf1_{\{X_i<x\}}\), with mean
\(F(x-)\), to obtain \(F_n(x-)\to F(x-)\) almost surely.
This supplies both kinds of endpoint control needed at atoms.

There are only countably many pairs of grid indices
\((j,k)\), \(k\geq2\), \(1\leq j<k\). Intersect the
probability-one events for both threshold conventions at every such
grid point. Call the resulting measurable event \(\Omega_1\).
At an outcome of \(\Omega_1\), for each fixed \(k\), the finitely
many convergences give an integer \(n_k\) such that for \(n\geq n_k\),
\[
 |F_n(x_{j,k})-F(x_{j,k})|\leq1/k,\qquad
 |F_n(x_{j,k}-)-F(x_{j,k}-)|\leq1/k
                         \quad(1\leq j<k).
\]
The infinite endpoints satisfy these estimates with error zero by
definition. If a real \(x\) equals any interior grid point, the
first estimate already bounds its error. Otherwise the finite
ordered grid, including its infinite endpoints, places \(x\) in
a nonempty open interval \((x_{j-1,k},x_{j,k})\).
Repeated grid points cause only empty intervals and introduce no
additional case.

Monotonicity, the endpoint estimates and the quantile-gap bound give
the complete upper comparison
\[
 \begin{aligned}
 F_n(x)&\leq F_n(x_{j,k}-)
       \leq F(x_{j,k}-)+1/k\\
       &\leq F(x_{j-1,k})+2/k
       \leq F(x)+2/k,
 \end{aligned}
\]
and the complete lower comparison
\[
 \begin{aligned}
 F_n(x)&\geq F_n(x_{j-1,k})
       \geq F(x_{j-1,k})-1/k\\
       &\geq F(x_{j,k}-)-2/k
       \geq F(x)-2/k.
 \end{aligned}
\]
These inequalities also apply to the two unbounded end intervals,
using the specified infinite-endpoint values. Thus \(D_n\leq2/k\)
for every \(n\geq n_k\), at the chosen outcome. This holds for
all \(k\geq2\) on \(\Omega_1\), so \(D_n\to0\) there.
Lemma~\ref{thm:p15-empirical-cdf-structure} already proved that \(D_n\) is measurable,
which makes the almost-sure assertion an assertion about measurable
random variables. The proof never assumes that \(F\) is continuous.
\end{proof}
\end{annexitem}

\begin{annexitem}{annex:p15.finite-alphabet-entropy}{thm:p15-finite-alphabet-entropy}{Finite product masses and the entropy inequalities}
\begin{proof}
For a fixed word \(w\), independence gives
\(\mathbb P(W_n=w)=\prod_{i=1}^np(w_i)=\pi_n(w)\).
Finite distributivity gives
\(\sum_{w\in\mathcal A^n}\pi_n(w)=(\sum_{a\in\mathcal A}p(a))^n=1\),
so these are normalized word masses, including the zero masses.
For each \(i\), \(\mathbb P(X_i\notin\mathcal A_+)=0\).
Their countable union is null. On its complement all factors in
every observed word probability are positive, and the logarithm of
a finite product is the sum of its logarithms. This proves the
information-sum identity simultaneously for every \(n\).

The real function \(I_0\) is measurable on the finite alphabet and
bounded, because the finitely many positive probabilities have a
positive minimum. Applying the same function to independent identically
distributed inputs gives i.i.d. real outputs. Direct finite summation
gives their mean \(H\) and their variance \(\sigma_I^2\).
Each supported probability is at most one, so \(-\log p(a)\geq0\)
and \(H\geq0\).

For \(t>0\), \(\log t\leq t-1\): the differentiable function
\(t-1-\log t\) decreases up to one, increases after one and is
zero there. Apply this inequality to \(t=1/(mp(a))\) at each
supported symbol and multiply by \(p(a)>0\). Summing yields
\[
 H-\log m
  =\sum_{a\in\mathcal A_+}p(a)\log\frac1{mp(a)}
  \leq\sum_{a\in\mathcal A_+}\left(\frac1m-p(a)\right)=0.
\]
This proves the upper entropy bound. When \(m=1\), its sole supported
symbol has probability one, so the information is identically zero
on the support and both \(H\) and \(\sigma_I^2\) are zero.
The value convention at a zero-mass symbol affects none of these
finite expectations or almost-sure identities.
\end{proof}
\end{annexitem}

\begin{annexitem}{annex:p15.finite-alphabet-aep}{thm:p15-finite-alphabet-aep}{The information average and limiting typical probability}
\begin{proof}
The preceding lemma makes \(I_0(X_i)\) a bounded i.i.d. real
sequence. The integrable strong law, Theorem~\ref{thm:p15-etemadi-strong-law}, gives
\(n^{-1}\sum_{i=1}^nI_0(X_i)\to H\) almost surely.
Intersect this event with the probability-one supported-sequence
event from the preceding lemma. On that intersection the information
sum equals \(-\log\pi_n(W_n)\) for every \(n\), proving the AEP.
Its definition with \(-\log0=\infty\) remains meaningful on the
measurable exceptional event as well.

For any word of positive probability, monotonicity of the logarithm
shows that the two non-strict typicality inequalities are equivalent to
\[
 \left|-\frac1n\log\pi_n(w)-H\right|\leq\varepsilon.
\]
A word of zero probability fails the positive lower bound and is
not typical. At any outcome where the information average converges
to \(H\), membership in the typical set is eventually true for
every fixed \(\varepsilon>0\). Its indicator therefore tends to one
almost surely and is bounded by one. Bounded convergence gives
\(\mathbb P(W_n\in\mathcal T_n(\varepsilon))\to1\).
If the support contains just one symbol, its observed word probability
is one and its information average and entropy are both zero; it
satisfies the same argument and inequalities.
\end{proof}
\end{annexitem}

\begin{annexitem}{annex:p15.typical-set-bounds}{thm:p15-typical-set-bounds}{The finite deviation bound and both counting inequalities}
\begin{proof}
Let \(L_n=n^{-1}\sum_{i=1}^nI_0(X_i)\).
Independence, finite variance addition and centering give
\(\mathbb EL_n=H\) and \(\operatorname{Var}(L_n)=\sigma_I^2/n\).
Outside a null event, atypicality is exactly \(|L_n-H|>\varepsilon\).
Chebyshev's inequality gives
\(\eta_n\leq\sigma_I^2/(n\varepsilon^2)\); the additional bound
\(\eta_n\leq1\) holds for every event. If \(\sigma_I^2=0\),
the same inequality says \(\eta_n=0\).

Sum the actual word probabilities over the finite typical set.
Each is at least \(e^{-n(H+\varepsilon)}\), so
\[
 |\mathcal T_n(\varepsilon)|e^{-n(H+\varepsilon)}
    \leq\sum_{w\in\mathcal T_n(\varepsilon)}\pi_n(w)
    =1-\eta_n\leq1.
\]
This gives the stated upper cardinality bound.
Each typical word also has probability at most \(e^{-n(H-\varepsilon)}\);
hence
\[
 1-\eta_n
   \leq|\mathcal T_n(\varepsilon)|e^{-n(H-\varepsilon)},
\]
which gives the lower cardinality bound. These inequalities remain
valid when the typical set is empty, or when \(H-\varepsilon<0\).

For any subset \(\mathcal C\) of the word space, split it into its
typical and atypical portions. The latter has probability at most
\(\eta_n\), and the former contains at most \(|\mathcal C|\)
words, each with probability at most \(e^{-n(H-\varepsilon)}\).
Therefore
\[
 \mathbb P(W_n\in\mathcal C)
       \leq\eta_n+|\mathcal C|e^{-n(H-\varepsilon)}.
\]
If its probability is at least \(1-\delta\), rearrangement gives
the stated lower bound on its cardinality, with zero replacing a
negative lower estimate. Finally the stated condition on \(n\)
makes the Chebyshev bound at most \(\delta\), giving typical
probability at least \(1-\delta\). All counting steps use the
joint word masses already established, rather than products of
marginals for an unspecified dependent process.
\end{proof}
\end{annexitem}

\begin{annexitem}{annex:p15.strong-law-application-calculations}{thm:p15-strong-law-application-calculations}{The renewal, atomic-distribution and word calculations}
\begin{proof}
For the zero-or-two waiting law the mean is
\(0(1/2)+2(1/2)=1\), and its probability of a positive waiting
time is \(1/2\). The renewal theorem therefore applies with
\(\mu=1\), giving limit one despite coincident epochs.
For the positive dyadic waiting law,
\(\sum_{j\geq1}2^{-j}=1\) and
\(\mathbb EX=\sum_{j\geq1}2^j2^{-j}=\infty\).
All values are positive finite reals. The same renewal theorem
therefore gives finite counts at finite times and limit zero for
their ratio to time. The required independent sequences exist by
the countable independent-sequence construction.

For the two-atom observation law, on one event of probability one all
observations lie in \(\{0,2\}\). At each such outcome,
\[
 \begin{array}{c|ccc}
 x&x<0&0\leq x<2&x\geq2\\ \hline
 F(x)&0&7/10&1\\
 F_n(x)&0&K_n/n&1
 \end{array}
\]
so the uniform error is precisely \(|K_n/n-7/10|\).
The quantile of each level \(j/10\leq7/10\) is zero, while
the quantiles of \(8/10\) and \(9/10\) are two. The displayed
grid therefore has exactly the repetitions stated. The ordinary
CDF value at zero is \(7/10\) whereas its left limit is zero,
which illustrates why both endpoint conventions were retained.

For the biased alphabet, the supported information values are
\[
 I_0(0)=\log4-\log3,\qquad I_0(1)=\log4.
\]
Their weighted mean is \(H=\log4-(3/4)\log3\). Their respective
deviations from this mean are \(-(1/4)\log3\) and
\((3/4)\log3\), so
\[
 \sigma_I^2
   =\frac34\frac1{16}(\log3)^2
       +\frac14\frac9{16}(\log3)^2
   =\frac3{16}(\log3)^2.
\]
Words containing \(\bot\) have probability zero. A word with \(k\)
ones and \(n-k\) zeros has probability
\((1/4)^k(3/4)^{n-k}=3^{n-k}4^{-n}\).
Its normalized negative logarithm is
\(\log4-(1-k/n)\log3\), whose difference from \(H\) is
\((k/n-1/4)\log3\). This proves the exact typicality condition.
There are \(\binom nk\) words with those counts, because their
\(k\) one-positions can be any \(k\)-element subset of the \(n\)
positions. Summing over the admissible integers proves the stated
cardinality; summing
\(\binom nk3^{n-k}4^{-n}\) over the same integers gives the exact
typical probability for any \(n,\varepsilon\).
At \(n=4\) and \(\varepsilon=(\log3)/8\), the allowed count satisfies
\(|k/4-1/4|\leq1/8\), so the only integer is \(k=1\).
There are four such words, each of probability \(27/256\), giving
total probability \(27/64\).

For the uniform law on its \(m\) supported symbols, every supported
information value is \(\log m\), so \(H=\log m\) and the variance
is zero. Every supported word has probability \(m^{-n}=e^{-nH}\),
which lies between the two typical bounds for any \(\varepsilon>0\).
No unsupported word belongs to the typical set. Its cardinality is
therefore \(m^n\), with all formulas still valid at \(m=1\).
\end{proof}
\end{annexitem}

\begin{annexitem}{annex:p15.strong-law-application-boundaries}{thm:p15-strong-law-application-boundaries}{The undefined-mean, uniformity and dependent-information failures}
\begin{proof}
For the symmetric dyadic law, the positive and negative masses sum
together to \(\sum_{j\geq1}2^{-j}=1\). Each part expectation is
\[
 \sum_{j\geq1}2^j2^{-j-1}=\sum_{j\geq1}1/2=\infty.
\]
For integers \(n\) in the block \(2^r\leq n<2^{r+1}\),
\[
 \mathbb P(|X_n|>n)=\sum_{j\geq r+1}2^{-j}=2^{-r}.
\]
There are \(2^r\) integers in this block, so each block contributes
one to the sum of these probabilities. The independent large-jump
events therefore occur infinitely often almost surely by the second
Borel--Cantelli lemma. On that event \(X_n/n\) does not tend to zero.
But a finite limit \(S_n/n\to L\) would imply, for \(n\geq2\),
\[
 X_n/n=S_n/n-\frac{n-1}{n}\frac{S_{n-1}}{n-1}\longrightarrow0.
\]
Thus finite convergence of the averages has probability zero.

It remains to exclude either infinite limit, rather than infer that
from the preceding argument. Write \(E_+=\{S_n/n\to+\infty\}\)
and \(E_-=\{S_n/n\to-\infty\}\). These are measurable events:
for example
\[
 E_+=\bigcap_{K\geq1}\bigcup_{N\geq1}
                      \bigcap_{n\geq N}\{S_n/n>K\}.
\]
Removing any fixed number of initial summands changes \(S_n/n\)
by a fixed finite outcome-dependent number divided by \(n\).
This does not change membership in \(E_+\) or \(E_-\), and
expressing the same events with tail sums shows that they belong
to every tail-coordinate sigma-algebra. The independent tail zero-one
law makes each probability zero or one.
The sequences \((X_n)\) and \((-X_n)\) have the same joint law:
symmetry and independence give equal probabilities on every finite
coordinate cylinder, and probability-measure uniqueness on that
generating pi-system extends equality to the product sigma-algebra.
Sign reversal exchanges the two countably defined events, so
\(\mathbb P(E_+)=\mathbb P(E_-)\). They are disjoint; their common
probability cannot be one and therefore is zero. Together with the
finite-limit exclusion, this proves absence of any extended-real
limit almost surely.

For all-zero waiting times, \(T_n=0\) at every index, and the
definition of \(N(t)\) includes all indices for \(t\geq0\).
For \(X_n=-1\), instead \(T_n=-n\leq t\) at all indices and
all nonnegative times. Both suprema are infinite. The latter example
violates nonnegativity and is not an admissible renewal model.

For the repeated fair observation \(X_i=B\in\{0,1\}\), the common
CDF and empirical CDF are both zero at \(x<0\), both one at
\(x\geq1\), and at \(0\leq x<1\) have values \(1/2\) and
\(\mathbf1_{\{B=0\}}\), respectively. Thus \(D_n=1/2\) for all
\(n\) at every outcome. The variables are dependent because the
events \(\{X_1=0\}\) and \(\{X_2=0\}\) coincide and have
probability \(1/2\), not joint probability \(1/4\).
For the deterministic CDFs \(G_n\), each is the distribution
function of unit mass at \(n\). For fixed real \(x\), it is
eventually zero, but \(G_n(n)=1\), giving uniform distance one
from the pointwise zero limit. That limit fails the CDF requirement
of tending to one at positive infinity.

For the all-Borel-set supremum, at a fixed outcome let
\(B_n=\{X_1,\ldots,X_n\}\). This is a finite Borel set with
\(\nu_n(B_n)=1\). The uniform law assigns zero mass to every
finite set, so \(\nu(B_n)=0\). The pointwise supremum over all
Borel sets is therefore at least one and at most one, since both
arguments are probabilities. The set selected here may depend on
the outcome; it is an allowable member of the class in each
pointwise supremum. This supplies the claimed failure of the
stronger uniformity assertion.

Finally take the single repeated biased bit. For the realized
constant word \(B^n\), the true joint probability is \(p(B)\),
while the marginal product is \(p(B)^n\); for \(n>1\) these
are different. The normalized negative logarithm of that product
is \(-\log p(B)\), at every \(n\).
Its difference from \(H\) is \(-(1/4)\log3\) if \(B=0\)
and \((3/4)\log3\) if \(B=1\). Both absolute differences exceed
\((\log3)/8\), so neither possible realized word is in the
product-law typical set at that tolerance. Its probability under
the dependent process is therefore zero for every \(n\).
Nevertheless the two words \(0^n,1^n\) have true probabilities
\(3/4,1/4\) and together exhaust the actual word law. This
also identifies the failed independence step: the product of
marginals does not provide the joint word masses on which the
typical-set estimates depend.
\end{proof}
\end{annexitem}
